\section{Results}

This section presents quantitative outcomes from
1 simulation scenario
tagged for analysis. Each scenario was optimized using the
forward beam search solver described in Section~II with beam
width $K = 100$. Figures and tables are generated directly from
the solver output; the accompanying narrative describes what each
chart shows and highlights the key patterns visible in the data.
Interpretation, trade-off analysis, and mission-planning
implications are addressed in Section~V (Discussion).


\subsection{Scenario:Currently Available Vehicles Minimizing Number of Launches}

Table~\ref{tab:metrics_1} summarizes the headline
performance metrics for this scenario.
The assembly campaign completed
58~modules in
154~periods
(35.9~months),
required 32~launch
events
(27~cargo and 5~crew
rotations),
and incurred a total program cost of
\$13660\,M.
The cumulative proximity risk metric
(Eq.~\ref{eq:risk}) accumulated to
0.0018 over the campaign.
These figures reflect a spacecraft of length
5\,km with
Chemical propulsion and a
Solar power system, optimized
under objective weights
$(w_1, w_2, w_3) =
(5,\;
 0,\;
 0)$.

\begin{table}[htbp]
\centering
\caption{Key performance metrics ---
         Currently Available Vehicles Minimizing Number of Launches.}
\label{tab:metrics_1}
\begin{tabular}{lr}
\toprule
Metric & Value \\
\midrule
Total launches        & 32 \\
\quad Cargo flights   & 27 \\
\quad Crew rotations  & 5 \\
Mission duration      & 154 periods
                        (35.9 months) \\
Total program cost    & \$13660\,M \\
Cost per module       & \$236\,M \\
Modules completed     & 58 \\
Cumulative risk       & 0.0018 \\
Spacecraft length     & 5\,km \\
Propulsion            & Chemical \\
Power system          & Solar \\
Cargo vehicles used &
  Falcon Heavy, Vulcan Centaur, SLS Block 2 \\
Crew vehicle used &
  Starliner \\
Objective weights ($w_1, w_2, w_3$) &
  (5,
   0,
   0) \\
\bottomrule
\end{tabular}
\end{table}

\paragraph{Assembly timeline (Fig.~\ref{fig:gantt_1}).}
The Gantt chart maps every action taken by the optimizer against
the period index (horizontal axis).
Each horizontal bar represents a single event: cargo launches
are shown in blue, crew launches in orange, fully assembled
modules in green, and work-in-progress (WIP) modules---those
whose assembly hours spanned more than one period---in purple.
Bars for concurrent actions within the same period are stacked
vertically, so a dense vertical column indicates a high-activity
period with simultaneous launches and assembly completions.
In this scenario the first cargo flight was dispatched at
period~4
(day~28,
or $\approx0.9$~months
after campaign start), which is gated by the lead time of the
selected cargo vehicles
(Eq.~\ref{eq:leadtime}).
The first crew rotation launched at period~13;
assembly work began at period~45, once
sufficient cargo had arrived and at least one module was
deliverable to the site.
The campaign logged 27~cargo
flights,
5~crew rotations,
and 58~module
completions
across 154~periods.

\begin{figure}[htbp]
  \centering
  \resizebox{0.85\textwidth}{!}{\input{generated/run_0_gantt.pgf}}
  \caption{Assembly timeline Gantt chart ---
           Currently Available Vehicles Minimizing Number of Launches.
           Blue~=~cargo launch; orange~=~crew launch;
           green~=~module assembled; purple~=~WIP.}
  \label{fig:gantt_1}
\end{figure}

\paragraph{Resource utilization (Fig.~\ref{fig:resources_1}).}
The resource chart overlays two time-varying quantities on a
dual-axis plot: the on-site crew headcount (left axis, solid
orange line) and the total vehicle count in proximity (right
axis, dashed blue line), where vehicle count includes all
docked crew vehicles plus any cargo vehicles actively unloading
in that period.
Together these two signals drive the congestion penalty
$\pi(n,p)$ in Eq.~(\ref{eq:penalty}): peak crew count
maximizes assembly throughput, but simultaneous peaks in both
crew and vehicle count incur the largest proximity penalties
and collision-risk contributions (Eq.~\ref{eq:risk}).
Periods in which the vehicle count spikes while crew count
remains constant correspond to cargo arrivals; those are the
highest-risk windows in the campaign.
Flat crew-count segments between rotation launches indicate
periods where crew is absent and assembly is stalled or purely
robotic.
The relatively elevated cumulative risk of
0.0018 in this scenario
is consistent with periods of high vehicle density visible in the
chart; tightening the crew-rotation overlap schedule or
staggering cargo arrivals would reduce the peak $n$ and
lower $\mathcal{R}$.

\begin{figure}[htbp]
  \centering
  \resizebox{0.75\textwidth}{!}{\input{generated/run_0_resources.pgf}}
  \caption{On-site crew (left axis, orange) and vehicle
           count in proximity (right axis, blue dashed) ---
           Currently Available Vehicles Minimizing Number of Launches.}
  \label{fig:resources_1}
\end{figure}

\paragraph{Program cost breakdown (Fig.~\ref{fig:cost_1}).}
The cost chart presents a two-panel horizontal bar chart.
The upper panel shows cargo launch vehicles (blue bars) and the
lower panel shows crew vehicles (orange bars).
Each bar's length equals the catalog unit price for that vehicle
multiplied by its total number of flights in this scenario; the
first-flight cost premium (Eq.~\ref{eq:firstflight}) is folded
into the first flight's bar.
Vehicles with zero flights are omitted; numeric cost labels
(\$M) appear at the right end of each bar.
The total program cost of \$13660\,M
is the sum of all bars in both panels.
Of the 32~total launches,
27~were cargo flights and 5~were
crew rotations; the relative height of the two panels indicates
the cargo-versus-crew share of the total budget.
A vehicle that appears in only one flight carries a
disproportionately high first-flight premium (up to 40\% for
in-development vehicles), which is visible as an extra bar
extension compared with per-launch cost alone.
Because $w_1 = 5$ is the
dominant weight, the optimizer consolidated payload onto
high-capacity vehicles to minimize the number of bars
in the upper panel.

\begin{figure}[htbp]
  \centering
  \resizebox{0.75\textwidth}{!}{\input{generated/run_0_cost.pgf}}
  \caption{Program cost breakdown by vehicle ---
           Currently Available Vehicles Minimizing Number of Launches.
           Upper panel: cargo vehicles (blue).
           Lower panel: crew vehicles (orange).
           Bar length = unit price $\times$ flights.}
  \label{fig:cost_1}
\end{figure}

\paragraph{Assembly progress (Fig.~\ref{fig:modules_1}).}
The modules-over-time chart is a step function of cumulative
modules completed plotted against period index.
Each vertical step represents one or more module completions
within the same period; the height of each step equals the
number of modules finished in that period.
A dotted horizontal red reference line marks the total module
target (58~modules); the campaign
is complete when the step function first reaches this line.
Flat horizontal segments---periods with no new completions---
indicate one of three conditions: (a) all eligible modules are
in transit and no material has yet arrived at the site;
(b) a WIP module is being worked across multiple periods and
will complete in a future step; or (c) the crew has expired
and no new rotation has yet arrived.
The overall slope of the curve is the assembly rate; a
consistently steep slope indicates that the crew and cargo
pipeline are well-matched, while a shallow slope or a late
inflection point suggests the campaign was bottlenecked by
transit time, crew availability, or DAG precedence constraints.
In this scenario, 58~modules were
completed in 154~periods, giving an
average assembly rate of
0.38~modules
per period
($\approx1.61$~per
calendar month).

\begin{figure}[htbp]
  \centering
  \resizebox{0.75\textwidth}{!}{\input{generated/run_0_modules.pgf}}
  \caption{Cumulative modules completed vs.\ period ---
           Currently Available Vehicles Minimizing Number of Launches.
           Dotted line = total module target
           (58~modules).}
  \label{fig:modules_1}
\end{figure}



\section{Discussion}
\label{sec:discussion}


\subsection{Solution Space and Constraint Boundaries}

Figure~\ref{fig:solution_space} maps the full set of solutions
generated by a $3 \times 3 \times 3$ grid of objective weights
$(w_1, w_2, w_3) \in \{0, 1, 2\}^3$ onto the
(launches, duration) plane.
Each point represents a unique (launches, duration) outcome
produced by the beam-search optimizer for one weight
combination; color indicates program cost.
Diamond markers identify Pareto-optimal solutions---points
for which no other sampled solution achieves a strictly
lower value on all three objectives---and the dashed step line
traces the empirical Pareto frontier through these points.
The red shaded band above the horizontal constraint line
($T = T_{\max}$) marks the infeasible region imposed by the
planning horizon limit; solutions approaching this boundary
indicate that the optimizer exhausted all available scheduling
periods without completing the spacecraft.
The orange band below the lead-time floor marks durations that
are physically unachievable regardless of weights, because
vehicle availability gates (Eq.~\ref{eq:leadtime}) prevent
first launches before the earliest $\tau_v$ has elapsed.
Red stars identify the tagged
scenario
chosen for detailed analysis above.
The density and spread of points in the interior of the plot
reveal how sensitive the campaign architecture is to weight
selection: a compact cluster indicates that most weight
combinations converge to similar solutions, while a widely
dispersed cloud shows strong weight sensitivity.

\begin{figure}[htbp]
  \centering
  \resizebox{0.85\textwidth}{!}{\input{generated/solution_space.pgf}}
  \caption{Complete solution space generated by a
           $3{\times}3{\times}3$ weight sweep.
           Diamond markers are Pareto-optimal points;
           the dashed line traces the empirical frontier.
           Red stars are the tagged
           scenario.
           Shaded bands mark the infeasible regions.}
  \label{fig:solution_space}
\end{figure}


The single tagged scenario
(Currently Available Vehicles Minimizing Number of Launches)
produced 32~launches,
154~periods of assembly
(35.9~months), and
\$13660\,M of total
program cost under weights
$(w_1, w_2, w_3) =
(5,\;
 0,\;
 0)$.
The discussion below interprets each performance dimension in light
of the model physics and the selected configuration.

\subsection{Launch Cadence and Vehicle Selection}

The 32-launch campaign includes
27~cargo
flights
and 5~crew
rotations,
averaging
0.21~total
launches per period
($\approx0.89$~per
calendar month).
The cargo fleet comprised
Falcon Heavy and Vulcan Centaur and SLS Block 2;
crew rotations used
Starliner.
At 0.55~launches per module,
the optimizer consolidated multiple modules onto individual
flights, reflecting the weight $w_1 = 5$
placed on minimizing launch count.
The relatively high per-flight payload utilization reduces
overall pad-time and range-coordination overhead.
The pad-turnaround constraint (Section~II.G) precluded any
single vehicle from flying more frequently than once every
$\Delta t_{\mathrm{pad}} = 2$ periods; the observed cadence
of 0.21~launches
per period is consistent with this ceiling for
32~flights across
154~periods.

\subsection{Schedule, Crew Coverage, and Assembly Bottlenecks}

The campaign completed in 154~periods
(35.9~months).
The first cargo flight launched at period~4,
gated by the lead time of the selected cargo vehicle.
Assembly work commenced at period~45,
after sufficient modules had arrived at the site and a crew
vehicle was docked to begin EVA operations.
The average assembly rate of
0.38~modules per period
(1.61~per calendar month) was
set by the interaction of crew EVA capacity
(Eqs.~\ref{eq:eva}--\ref{eq:crew_hours}), DAG precedence
constraints, and proximity congestion penalties
(Eq.~\ref{eq:penalty}).
The dominant launch weight $w_1 = 5$
caused the optimizer to defer some launches, accepting a longer
schedule in exchange for higher per-flight payload utilization.
Flat regions in the modules-over-time curve
(Fig.~\ref{fig:modules_1}) correspond to these deliberate
deferral windows.
Gaps between crew rotations---periods in which the crew vehicle
had expired and no replacement had yet arrived---are the primary
source of schedule inefficiency; their aggregate contribution
to the campaign length is visible as flat segments in
Fig.~\ref{fig:modules_1}.

\subsection{Cost Structure and Program Economics}

At \$13660\,M total,
the campaign incurred an average cost of
\$236\,M per
assembled module.
The cost breakdown (Fig.~\ref{fig:cost_1}) decomposes this
total into contributions from each cargo and crew vehicle.
Because cost was not the dominant objective
($w_3 = 0$), the optimizer may have
selected higher-capacity or faster-available vehicles at a cost
premium; the cost chart shows where the largest single-vehicle
expenditures occurred.
Crew rotations typically represent a significant share of total
program cost because each rotation incurs both a launch fee and
a crew-support overhead; reducing rotation frequency (via longer
crew tours) reduces cost but may leave the assembly site
uncrewed during transit windows.

\subsection{Proximity Risk and Operational Safety}

The cumulative proximity risk of
0.0018 was accumulated
over 154~periods according to
Eq.~(\ref{eq:risk}).
The low cumulative risk reflects well-separated vehicle traffic:
cargo and crew arrivals were largely non-overlapping, keeping
the simultaneous vehicle count $n(t)$ close to~1 for most of the
campaign and the $\binom{n}{2}$ pair contribution near zero.
This favorable safety profile was achieved without an explicit
safety objective, suggesting that the campaign structure
naturally minimizes multi-vehicle overlap when launch timing is
driven by assembly-demand signals.

